\documentclass[11pt,a4paper]{article}

\usepackage{setspace}
\usepackage{geometry}
\usepackage{titlesec}
\usepackage{array}
\usepackage{longtable}
\usepackage{graphicx}
\usepackage{caption}
\usepackage{float}
\usepackage{titling}
\usepackage{enumitem}
% --- Tighter margins for 2-page fit ---
\geometry{
 a4paper,
 left=0.9in,
 right=0.9in,
 top=0.9in,
 bottom=0.9in
}
\setlength{\droptitle}{-9em}
% --- Tighter spacing ---
\setstretch{1.12}

% --- Reduce spacing before/after section titles ---
\titlespacing*{\section}{0pt}{4pt}{4pt}

% --- Reduce paragraph spacing ---
\setlength{\parskip}{6pt}

\title{\textbf{Abstract} \\ \large \textbf{College Bus Surveillance using Computer Vision}}
\author{}
\date{}

\begin{document}

\maketitle
\vspace{-3.5cm}
\section*{a) Problem Statement}
Ensuring student safety and accurate bus attendance remains a challenge for institutions. Manual roll-calls and paper logs are slow, error-prone, and lack real-time visibility. RFID/card-tap systems can fail when students forget cards and involve additional hardware, making them unreliable and costly. A fully automated, low-cost, real-time monitoring solution is required.

\section*{b) Solution Approach}
The proposed system offers an intelligent college bus monitoring solution using mobile-based computer vision, GPS tracking, and centralized cloud storage. No specialized hardware such as Raspberry Pi or edge devices is required; instead, modern smartphones serve as the processing unit. The system includes:  
1) Centralized student database and web dashboard,  
2) Master mobile app for face recognition,  
3) Real-time GPS logging and attendance tracking.

\section*{c) Technologies Used}
Frontend: Flutter. Backend: Python Flask.  
Firestore stores student data, embeddings, GPS logs, and attendance.  
Face detection uses Google ML Kit; recognition uses TensorFlow Lite.  
Additional integrations: camera processing, Google Maps SDK, Directions API, Geolocation API, and OpenStreetMap for route visualization.

\section*{d) Base Paper Details}

\subsection*{Base Paper 1}
\textbf{“Bus Attendance System Through Facial Recognition by Implementing AI Integration”}  
Authors: Johnson C, Karthick M, Sakthivel P, Sasikala K, Santhi A  
Published: 2023, \textit{IJRPR}, Vol. 8, Issue 11  
Link: \href{https://www.ijnrd.org/papers/IJNRD2311385.pdf}{IJRPR}

\subsection*{Base Paper 2}
\textbf{“Smart College Bus Tracker: Real-Time Location and ETA Predictor”}  
Authors: M. Sowndharya, Yuvaraj K, Rahul S, Srichandiran S, Sriram S  
Published: 2025, \textit{IJSRED}, Vol. 8, Issue 3  
DOI: IJSRED-V8I3P68  
Link: \href{https://ijsred.com/volume8/issue3/IJSRED-V8I3P68.pdf}{IJSRED}

\section*{e) Group No. \& Team Members}
\textbf{Group No: 6}
\begin{itemize}
    \item Akash V G – 8
    \item Fidha Fathima A – 32
    \item Mohammed Ramshad T S – 44
    \item Nithya Manoj – 47
\end{itemize}

\section*{f) Time Schedule}

% smaller font inside longtable to save space
{\small
\begin{longtable}{|c|p{10cm}|}
\hline
\textbf{Week} & \textbf{Work to be Completed} \\ \hline

Week 1 & Requirement analysis, scope finalization, architecture diagram, module identification, UI wireframing. \\ \hline

Week 2 & DFD/ER diagrams, Flask backend setup, Firestore structure, authentication, REST API development. \\ \hline

Week 3 & Web dashboard: admin login, student management, bus assignment, attendance logs, map setup. \\ \hline

Week 4 & Flutter app: camera preview, ML Kit detection, UI (login, camera, bus selection). \\ \hline

Week 5 & TensorFlow Lite recognition: embedding extraction, comparison, threshold tuning, accuracy tests. \\ \hline

Week 6 & Attendance workflow: GPS + timestamp logs, duplicate prevention, offline sync. \\ \hline

Week 7 & Real-time GPS tracking, Firestore updates, route plotting, Google Maps/OSM integration. \\ \hline

Week 8 & System integration and testing, low-light evaluation, network failure handling, optimization. \\ \hline

Week 9 & Security, UI polishing, performance tuning, admin filters, export features. \\ \hline

Week 10 & LaTeX report preparation. \\ \hline
\end{longtable}
}
\end{document}
